\section*{Chapter \ref{ch:g12:destruction-of-mountains} --- The Destruction of Mountains}

\begin{solution}{exo:g12:destruction-of-mountains:1}
(a) the \omterm{def:g12:destruction-of-mountains:denudation}{denudation rate}; (b) the \omterm{def:g12:destruction-of-mountains:load}{bed load}; (c) the \omterm{def:g12:destruction-of-mountains:load}{dissolved load};
(d) \omterm{def:g12:destruction-of-mountains:diagenesis}{diagenesis}; (e) \omterm{def:g12:destruction-of-mountains:rebound}{isostatic rebound}.
\end{solution}

\begin{solution}{exo:g12:destruction-of-mountains:2}
(1) A conglomerate. (2) A sandstone. (3) A shale. (4) A breccia.
\end{solution}

\begin{solution}{exo:g12:destruction-of-mountains:3}
\qty{1}{m/Myr} is \qty{1000}{mm} in \num{1000000} years, \qty{0.001}{mm} a
year. (1) \qty{0.054}{mm} a year. (2) \qty{0.218}{mm} a year.
(3) $\qty{10}{mm}$ a year is \qty{10}{m} in 1000 years, \qty{10000}{m/Myr}.
\end{solution}

\begin{solution}{exo:g12:destruction-of-mountains:4}
The solid load is \qty{85}{\percent} of the total, so the total is $60/0.85
\approx 70.6$ million tonnes, and the \omterm{def:g12:destruction-of-mountains:load}{dissolved load} about $70.6 - 60 \approx
10.6$ million tonnes a year.
\end{solution}

\begin{solution}{exo:g12:destruction-of-mountains:5}
Weathering of the granite, transport by a river, deposition in a basin,
burial under younger layers, cementation.
\end{solution}

\begin{solution}{exo:g12:destruction-of-mountains:6}
At \qty{30}{cm/s} the sand grain is picked up (above \qty{20}{cm/s}); the
silt grain stays in suspension. At \qty{5}{cm/s} the sand grain settles
(below \qty{10}{cm/s}), but the silt grain still stays in suspension, since
\qty{0.1}{cm/s} is enough to keep it there.
\end{solution}

\begin{solution}{exo:g12:destruction-of-mountains:7}
Total load $4/0.85 \approx 4.71$ million tonnes, $\num{4.71e9}/2830 \approx
\qty{1.66e6}{m^3}$ of rock a year. Over $\qty{5e10}{m^2}$: $\num{3.3e-5}$ m a
year, about \qty{33}{m/Myr}, below the median of \qty{54}{m/Myr}.
\end{solution}

\begin{solution}{exo:g12:destruction-of-mountains:8}
(1) $1 - 2700/3300 \approx 0.182$. (2) About \qty{182}{m}. (3) A little more
than the \qty{163}{m} of the chapter: the lighter the crust compared with the
mantle, the less the column rises.
\end{solution}

\begin{solution}{exo:g12:destruction-of-mountains:9}
(1) \qty{8}{km} of rock. (2) $8000 \times 0.163 \approx \qty{1300}{m}$.
(3) $8000/200 = 40$ million years.
\end{solution}

\begin{solution}{exo:g12:destruction-of-mountains:10}
(1) $794/43 \approx 18$ million years. (2) The surface goes down by only
$0.163$ times the thickness eroded, so the rock to remove is $794/0.163
\approx \qty{4900}{m}$, which takes about $4900/43 \approx 110$ million years.
(3) The continents are not only worn down: mountain building and the
magmatism of arcs keep raising them and adding new crust.
\end{solution}

\begin{solution}{exo:g12:destruction-of-mountains:11}
(1) $\num{1e14}/2830 \approx \qty{3.5e10}{m^3}$, that is about
\qty{35}{km^3}. (2) $100/15.5 \approx 6.5$ years.
\end{solution}

\begin{solution}{exo:g12:destruction-of-mountains:12}
$E = kh$ and $h = h_0 e^{-fkt}$ give $E(t) = kh_0 e^{-fkt}$. The thickness
removed from $t = 0$ to $t \to \infty$ is $\int_0^{\infty} kh_0
e^{-fkt}\,\mathrm{d}t = kh_0/(fk) = h_0/f$. To lower the surface by $h_0$,
erosion must remove $h_0/f \approx 6.1\,h_0$, as the rebound proposition
says: each metre eroded lowers the surface by only $f$ metres.
\end{solution}

\begin{solution}{exo:g12:destruction-of-mountains:13}
(1) $10 \times 1000 = \qty{10000}{mm}$, \qty{10}{m}. (2) Removing a load of
ice of thickness $t$ is like removing a mass $\rho_{\text{ice}}\,t$ per unit
area; the column rises until it has displaced that much less mantle:
$2000 \times 900/3380 \approx \qty{530}{m}$. (3) The land rises less and
less as it nears its new equilibrium: what drives the rise is the part not
yet recovered.
\end{solution}

\begin{solution}{exo:g12:destruction-of-mountains:14}
(1) Before 1900: about \qty{56}{m/Myr} (as in the example). In
1987--2006: $145/0.85 \approx 171$ million tonnes, $\num{1.71e11}/2830
\approx \qty{6.0e7}{m^3}$, over $\qty{2.98e12}{m^2}$: about
\qty{20}{m/Myr}. (2) The earlier one: the later load is held back by dams
and other works. (3) No: the decrease is $400 - 145 = 255$ million tonnes a
year, of which the dams account for 100 to 150, about half.
\end{solution}

\begin{solution}{exo:g12:destruction-of-mountains:15}
(1) At least \qty{6}{km}. (2) $310 - 252 = 58$ million years at most, so at
least $6000/58 \approx \qty{100}{m/Myr}$. (3) It is between the median
(\qty{54}{m/Myr}) and the mean (\qty{218}{m/Myr}) of the basins: an ordinary
rate for a mountain belt.
\end{solution}

\begin{solution}{pb:g12:destruction-of-mountains:1}
\textbf{1.} $E = 0.209 \times 2000 \approx \qty{418}{m/Myr}$.

\textbf{2.} \qty{0.418}{mm} a year.

\textbf{3.} Faster than both: it is a steep, young range, as fast as the
seismically active basins and more.

\textbf{4.} Its mean slope.

\textbf{5.} A layer of $\qty{0.418e-3}{m}$ over $\qty{1e6}{m^2}$ is
\qty{418}{m^3}, $418 \times 2830 \approx \qty{1.18e6}{kg}$: about 1200
tonnes.

\textbf{6.} $2000/418 \approx 4.8$ million years.

\textbf{7.} $h(t) = 2000\,e^{-0.209\,t}$ ($t$ in millions of years).

\textbf{8.} $1/0.209 \approx 4.8$ million years.

\textbf{9.} $\ln 2/0.209 \approx 3.3$ million years.

\textbf{10.} $e^{-0.209 \times 11} = e^{-2.30} \approx 0.100$.

\textbf{11.} $1 - 2830/3380 \approx 0.163$.

\textbf{12.} The surface goes down by $100 \times 0.163 \approx \qty{16}{m}$;
the \omterm{def:g12:continental-crust:root}{root} rises by about \qty{84}{m}.

\textbf{13.} $h' = -0.163 \times 0.209\,h \approx -0.0341\,h$, so $h(t) =
2000\,e^{-0.0341\,t}$.

\textbf{14.} $(2000 - 200)/0.163 \approx \qty{11000}{m}$ of rock.

\textbf{15.} Rocks that were about \qty{11}{km} down: metamorphic rocks
such as schists and gneisses, and granites.

\textbf{16.} $1/0.0341 \approx 29$ million years.

\textbf{17.} $29 \times \ln 10 \approx 67$ million years.

\textbf{18.} $2000\,e^{-0.0341 \times 200} \approx 2000 \times 0.0011 \approx
\qty{2}{m}$.

\textbf{19.} $2000\,e^{-0.0341 \times 250} \approx 2000 \times 0.0002 \approx
\qty{0.4}{m}$: a flat plain. A summit of \qty{2037}{m} means that the model
alone does not fit: the region must have been raised again since, or eroded
more slowly than the model assumes.

\textbf{20.} The mean elevation halves in $29 \times \ln 2 \approx$
\textbf{20 million years}.
\end{solution}
